\documentclass[../cours_2de.tex]{subfiles}
\begin{document}
\setcounter{chapter}{22}
\chapter{Algorithmique et Python}
\label{sec:ch23}
\begin{objectif}
Lire, écrire et tester un programme ; comprendre variables, fonctions et boucles ; relier chaque algorithme à un problème mathématique. Ce chapitre est un atelier de référence à utiliser tout au long du cours.
\end{objectif}
\section{Variables, types et affectation}
Une variable informatique mémorise une valeur. Les types rencontrés sont
\texttt{int} (entier), \texttt{float} (réel approché), \texttt{bool} (vrai ou faux), \texttt{str} (texte).
L'affectation en langage naturel $x\leftarrow x+1$ se note \texttt{x = x + 1} : on calcule d'abord avec l'ancienne valeur, puis on la remplace.
Le test d'égalité se note \texttt{==}, la différence \texttt{!=}.
\begin{lstlisting}[language=Python]
x = 4
x = x + 1
est_pair = (x % 2 == 0)
message = "La valeur obtenue est"
print(message, x, est_pair)
\end{lstlisting}
Ce programme affiche une valeur $5$ et un booléen faux. Les instructions sont exécutées dans l'ordre.
Avec deux variables, échanger leurs valeurs demande une variable temporaire ou l'affectation simultanée \texttt{x, y = y, x}.
Les nombres flottants sont approchés : \texttt{0.1 + 0.2 == 0.3} peut être faux. Un test numérique utilise une tolérance justifiée ; un calcul exact peut utiliser des entiers ou des fractions.
\section{Conditions et fonctions}
\begin{lstlisting}[language=Python]
def tarif(age):
    if age < 12:
        return 6
    elif age < 18:
        return 9
    else:
        return 12

def distance(ax, ay, bx, by):
    return ((bx-ax)**2 + (by-ay)**2)**0.5
\end{lstlisting}
\texttt{def} définit une fonction, ses arguments reçoivent les valeurs lors de l'appel, et \texttt{return} renvoie le résultat.
\texttt{tarif(15)} renvoie $9$ ; \texttt{distance(0,0,3,4)} renvoie $5$.
Une fonction informatique peut avoir plusieurs arguments ; une fonction mathématique réelle à une variable est un cas particulier.
L'indentation définit les blocs. \texttt{print} affiche un résultat mais ne remplace pas \texttt{return}.
\section{Boucles bornées et non bornées}
\begin{lstlisting}[language=Python]
def somme_impairs(n):
    total = 0
    for k in range(n):
        total = total + (2*k + 1)
    return total

def seuil(q, cible):
    # q > 1, cible >= 1
    n, valeur = 0, 1
    while valeur <= cible:
        n, valeur = n + 1, valeur * q
    return n
\end{lstlisting}
\texttt{range(n)} fournit les entiers de $0$ à $n-1$.
La boucle \texttt{for} a un nombre de tours connu ; la boucle \texttt{while} continue tant que la condition est vraie.
Avant chaque nouveau tour de la seconde boucle, \texttt{valeur} représente $q^n$ : c'est un \emph{invariant}. La preuve de terminaison exige $q>1$, comme au chapitre 4.
\begin{attention}
Une boucle \texttt{while} qui ne modifie pas la variable de son test peut tourner indéfiniment. Écrire d'abord les conditions d'entrée, l'état après un tour et la condition d'arrêt.
\end{attention}
\section{Un balayage exact pour encadrer une racine}
Le programme suivant renvoie les entiers $k,k+1,d$ tels que
$k/d\leq\sqrt2<(k+1)/d$, avec $d=10^n$.
\begin{lstlisting}[language=Python]
def encadrement_racine2(n):
    k, d = 1, 1
    for _ in range(n):
        k, d = 10*k, 10*d
        while (k+1)**2 <= 2*d*d:
            k = k + 1
    return k, k+1, d
\end{lstlisting}
On cherche chiffre après chiffre le dernier entier dont le carré ne dépasse pas $2d^2$.
Les comparaisons se font sur des entiers, ce qui évite l'ambiguïté d'un arrondi.
\texttt{encadrement\_racine2(3)} renvoie $(1414,1415,1000)$.
\section{Dichotomie et longueur polygonale}
\begin{lstlisting}[language=Python]
def dichotomie_racine2(n):
    a, b = 1.0, 2.0
    for _ in range(n):
        m = (a+b)/2
        if m*m < 2:
            a = m
        else:
            b = m
    return a, b
\end{lstlisting}
L'amplitude théorique est $2^{-n}$, jusqu'aux limites de précision des flottants. Pour un nombre excessif de tours, les deux bornes finissent par ne plus être distinguées par la machine.
\begin{lstlisting}[language=Python]
def longueur_polygone(f, a, b, n):
    pas = (b-a)/n
    total = 0.0
    for k in range(n):
        x, y = a+k*pas, a+(k+1)*pas
        total += (pas**2 + (f(y)-f(x))**2)**0.5
    return total
\end{lstlisting}
On exige $a<b$ et un entier $n\geq1$.
Pour la droite $f(x)=2x$ entre $0$ et $1$, chaque subdivision donne exactement $\sqrt5$ en arithmétique réelle, car la ligne polygonale est déjà le segment. Pour une parabole, comparer plusieurs subdivisions illustre l'approximation.
\section{Moyenne et écart type}
Le paramètre \texttt{valeurs} représente une collection de nombres ; les détails sur les listes sont un support, pas un prérequis obligatoire.
\begin{lstlisting}[language=Python]
from math import sqrt

def moyenne_ecart_type(valeurs):
    n = len(valeurs)
    if n == 0:
        raise ValueError("Serie vide")
    m = sum(valeurs) / n
    v = sum((x-m)**2 for x in valeurs) / n
    return m, sqrt(v)
\end{lstlisting}
Lire le programme consiste à identifier l'effectif, la moyenne, la moyenne des écarts quadratiques et sa racine. On utilise la variance descriptive de dénominateur $n$.
\section{Simuler des répétitions}
\begin{lstlisting}[language=Python]
from random import random

def experience(p):
    return int(random() < p)

def frequence(p, n):
    succes = 0
    for _ in range(n):
        succes += experience(p)
    return succes / n
\end{lstlisting}
On impose $0\leq p\leq1$ et $n\geq1$. Le générateur simule un uniforme sur $[0,1[$ ; la condition \texttt{random() < p} modélise un succès de probabilité $p$.
Répéter \texttt{frequence(0.5,100)} donne des résultats différents : c'est la fluctuation d'échantillonnage. En augmentant $n$, les grands écarts deviennent moins fréquents, sans imposer une amélioration à chaque exécution.
La graine fixe une suite reproductible pour un essai ; elle ne rend pas le modèle plus vrai.
\section{Filtrer et tabuler}
Les connecteurs Python \texttt{and}, \texttt{or}, \texttt{not} traduisent et, ou, non. Pour deux listes représentant des caractères des mêmes individus, parcourir simultanément leurs éléments, tester le critère et augmenter les compteurs.
Le dossier \texttt{python/} contient des versions exécutables avec contrôles d'entrée : divisibilité, géométrie, recherche de seuil, encadrements, balayage, statistiques, simulations et lecture du fichier cantonal. Le script de vérification teste des cas connus et des propriétés mathématiques.
\utilite{Programmer oblige à rendre explicites les hypothèses et l'ordre du raisonnement. Un programme correct sur un exemple doit encore être confronté à des cas limites et à une justification générale.}
\begin{retenir}
Décrire le problème avant le code, tracer les variables, tester une réponse connue, traiter les cas limites et séparer simulation, approximation et démonstration.
\end{retenir}
\end{document}
